\begin{theorem}[对角作用下的 \texorpdfstring{$q$}{q}-元轨道分类与受限 Bell 计数]\label{thm:op-algebra-fold-gauge-orbit-classification-qtuple}
\leanverified{paper\_op\_algebra\_fold\_gauge\_orbit\_classification\_qtuple}
令 \(q\ge 1\)。
令 \(G_m\) 按对角方式作用于 \(\Omega_m^q\)：
\(
\sigma\cdot(a_1,\dots,a_q):=(\sigma(a_1),\dots,\sigma(a_q))
\)。
对 \(\mathbf a=(a_1,\dots,a_q)\in\Omega_m^q\) 记
\[
x_i:=\Fold_m(a_i)\in X_m,
\qquad
I_x(\mathbf a):=\{i\in\{1,\dots,q\}: x_i=x\}.
\]
并在每个 \(I_x(\mathbf a)\) 上定义等价关系
\[
i\sim_x j\quad\Longleftrightarrow\quad a_i=a_j,
\]
从而得到集合划分 \(\Pi_x(\mathbf a)\)（其块数不超过 \(d_m(x)\)）。
则对任意 \(\mathbf a,\mathbf b\in\Omega_m^q\) 有
\[
\mathbf b\in G_m\cdot \mathbf a
\quad\Longleftrightarrow\quad
\bigl(\Fold_m(a_i)\bigr)_{i=1}^q=\bigl(\Fold_m(b_i)\bigr)_{i=1}^q
\ \text{且}\
\Pi_x(\mathbf a)=\Pi_x(\mathbf b)\ \ (\forall x\in X_m).
\]
因此轨道空间 \(\Omega_m^q/G_m\) 与“以 \(X_m\) 为颜色、每色的块数受限于 \(d_m(x)\) 的彩色集合划分族”等势。
其轨道数满足闭式
\[
\bigl|\Omega_m^q/G_m\bigr|
=
\sum_{\substack{(n_x)_{x\in X_m}\in\mathbb Z_{\ge 0}^{X_m}\\ \sum_x n_x=q}}
\frac{q!}{\prod_x n_x!}\;
\prod_{x\in X_m} B_{n_x}^{(\le d_m(x))},
\]
其中
\(
B_n^{(\le d)}:=\sum_{k=0}^{\min(n,d)} S(n,k)
\)
为“至多 \(d\) 块”的受限 Bell 数，而 \(S(n,k)\) 为第二类 Stirling 数。
\end{theorem}

\begin{proof}
\(\Rightarrow\) 方向：若 \(\mathbf b=\sigma\cdot\mathbf a\) 且 \(\sigma\in G_m\)，则 \(\Fold_m\circ\sigma=\Fold_m\)，从而两者的 \((x_i)\) 相同。
并且 \(\sigma\) 在每条纤维 \(\Fold_m^{-1}(x)\) 上为双射，因而保持“相等/不等”关系，故每个 \(x\) 的划分 \(\Pi_x\) 不变。

\(\Leftarrow\) 方向：若两者具有相同的 \((x_i)\) 与 \(\Pi_x\)，则对每个 \(x\)，两 \(q\)-元组在纤维 \(\Fold_m^{-1}(x)\) 中使用了同样数目的互异微态，并且重复模式一致。
由于 \(|\Fold_m^{-1}(x)|=d_m(x)\) 足以容纳该块数，存在某个 \(\sigma_x\in\mathfrak S_{d_m(x)}\) 把 \(\mathbf a\) 在该纤维中出现的互异微态逐块送到 \(\mathbf b\) 对应微态。
将各 \(\sigma_x\) 在不交并纤维上拼接得到 \(\sigma\in\prod_x\mathfrak S_{d_m(x)}\cong G_m\)，使 \(\sigma\cdot\mathbf a=\mathbf b\)。

计数式：先按颜色计数 \(n_x:=|I_x|\) 分组，分组方式为多项式系数 \(q!/\prod_x n_x!\)。
在固定 \(n_x\) 下，纤维内的重复模式等价于 \(I_x\) 的一个集合划分，其块数至多为 \(d_m(x)\)，计数即 \(B_{n_x}^{(\le d_m(x))}\)。
对颜色取乘积并对 \((n_x)\) 求和即得。证毕。
\end{proof}

\begin{theorem}[Burnside 轨道矩与纤维分解的指数矩母函数]\label{thm:op-algebra-fold-orbit-mgf-fiber-factorization}
\leanverified{paper\_op\_algebra\_fold\_orbit\_mgf\_fiber\_factorization}
对整数 \(q\ge 0\) 记
\[
\mathrm{Orb}_m(q):=\bigl|\Omega_m^q/G_m\bigr|.
\]
令 \(\sigma\) 在有限群 \(G_m\) 上服从均匀分布，并记其在 \(\Omega_m\) 上的不动点数
\[
F(\sigma):=\bigl|\mathrm{Fix}_{\Omega_m}(\sigma)\bigr|.
\]
则对任意 \(q\ge 0\) 有 Burnside 恒等式
\[
\boxed{\;\mathrm{Orb}_m(q)=\mathbb E\!\left[F(\sigma)^q\right]\;}
\]
并且指数矩母函数满足严格分解
\[
\sum_{q\ge 0}\mathrm{Orb}_m(q)\frac{z^q}{q!}
=\mathbb E[e^{zF(\sigma)}]
=\prod_{x\in X_m}\left(\sum_{k=0}^{d_m(x)}\frac{(e^{z}-1)^k}{k!}\right).
\]
\end{theorem}
\begin{proof}
对角作用的 Burnside 引理给出
\[
\mathrm{Orb}_m(q)
=\frac{1}{|G_m|}\sum_{\sigma\in G_m}\bigl|\mathrm{Fix}_{\Omega_m^q}(\sigma)\bigr|
=\frac{1}{|G_m|}\sum_{\sigma\in G_m}\bigl|\mathrm{Fix}_{\Omega_m}(\sigma)\bigr|^q
=\mathbb E[F(\sigma)^q],
\]
其中第二步使用 \(\mathrm{Fix}_{\Omega_m^q}(\sigma)=\mathrm{Fix}_{\Omega_m}(\sigma)^q\)。
对任意 \(z\) 求指数生成函数即得第一条等式。
\par
由 \(G_m\cong\prod_{x\in X_m}\mathfrak S_{d_m(x)}\)，可将 \(\sigma\) 写为 \(\sigma=(\sigma_x)_{x\in X_m}\) 且各 \(\sigma_x\in\mathfrak S_{d_m(x)}\) 独立均匀。
此时 \(F(\sigma)=\sum_{x\in X_m}\mathrm{Fix}(\sigma_x)\)，从而
\(
\mathbb E[e^{zF(\sigma)}]=\prod_x \mathbb E[e^{z\,\mathrm{Fix}(\sigma_x)}]
\)。
对 \(d\ge 0\) 令 \(\pi\) 为 \(\mathfrak S_d\) 的均匀随机置换。
对任意整数 \(0\le k\le d\)，其下降阶乘矩满足 \(\mathbb E[(\mathrm{Fix}(\pi))_k]=1\)：因为任取 \(k\) 个不同点同时被固定的概率为 \((d-k)!/d!=1/(d)_k\)，而 \((\mathrm{Fix})_k\) 计数有序的 \(k\)-元不动点组。
由恒等式
\(
(1+t)^{\mathrm{Fix}}=\sum_{k\ge 0}(\mathrm{Fix})_k\,t^k/k!
\)
并取期望，得 \(\mathbb E[(1+t)^{\mathrm{Fix}(\pi)}]=\sum_{k=0}^{d}t^k/k!\)。
令 \(t=e^{z}-1\) 即得
\(
\mathbb E[e^{z\,\mathrm{Fix}(\pi)}]=\sum_{k=0}^{d}(e^{z}-1)^k/k!
\)。
代回并对 \(x\) 取乘积得到结论。证毕。
\end{proof}

\begin{theorem}[低阶轨道的 Touchard 通用律]\label{thm:op-algebra-fold-orbit-touchard-loworder}
记最小纤维多重度
\[
d_{\min}(m):=\min_{x\in X_m} d_m(x).
\]
则对一切整数 \(1\le q\le d_{\min}(m)\) 有严格闭式
\[
\boxed{\;\mathrm{Orb}_m(q)=T_q(|X_m|)\;}
\]
其中 Touchard 多项式定义为 \(T_q(\lambda):=\sum_{k=0}^{q}S(q,k)\lambda^k\)，\(S(q,k)\) 为第二类 Stirling 数。
\end{theorem}
\leanverified{paper\_op\_algebra\_fold\_orbit\_touchard\_loworder}
\begin{proof}
由定理 \ref{thm:op-algebra-fold-gauge-orbit-classification-qtuple} 的计数式，当 \(q\le d_{\min}(m)\) 时对任意满足 \(\sum_x n_x=q\) 的向量都有 \(n_x\le d_m(x)\)，从而
\(
B_{n_x}^{(\le d_m(x))}=\sum_{k=0}^{n_x}S(n_x,k)=:B_{n_x}
\)
为普通 Bell 数。
因此指数生成函数化简为
\[
\sum_{q\ge 0}\mathrm{Orb}_m(q)\frac{z^q}{q!}
=\left(\sum_{n\ge 0}B_n\frac{z^n}{n!}\right)^{|X_m|}
=\exp\bigl(|X_m|\,(e^{z}-1)\bigr),
\]
其中使用了经典恒等式 \(\sum_{n\ge 0}B_n z^n/n!=\exp(e^{z}-1)\)。
右端的 \(q\) 阶导数在 \(z=0\) 处即为 Touchard 多项式 \(T_q(|X_m|)\)，从而得到结论。证毕。
\end{proof}

\leanverified{paper\_op\_algebra\_fold\_orbit\_first\_deviation\_minsector}
\begin{proposition}[首次偏离项与最小退化扇区规模]\label{prop:op-algebra-fold-orbit-first-deviation-minsector}
令 \(q_0:=d_{\min}(m)+1\)，并记最小退化扇区规模
\[
N_{\min}(m):=\#\{x\in X_m:\ d_m(x)=d_{\min}(m)\}.
\]
则有精确恒等式
\[
\boxed{\;\mathrm{Orb}_m(q_0)=T_{q_0}(|X_m|)-N_{\min}(m)\;}
\]
从而 \(N_{\min}(m)=T_{q_0}(|X_m|)-\mathrm{Orb}_m(q_0)\) 可由一步轨道计数完全确定。
\end{proposition}
\begin{proof}
在定理 \ref{thm:op-algebra-fold-gauge-orbit-classification-qtuple} 的求和式中，当 \(q=q_0\) 时仅有一种情形会触发受限：存在某个 \(x\) 使 \(n_x=q_0\) 且 \(d_m(x)=d_{\min}(m)\)。
此时
\(
B_{q_0}^{(\le d_{\min}(m))}=B_{q_0}-1
\)
（缺少唯一的 \(q_0\) 块划分）。
每个满足 \(d_m(x)=d_{\min}(m)\) 的颜色对该项的修正贡献恰为 \(-1\)，故总修正为 \(-N_{\min}(m)\)；其余项与 \(T_{q_0}(|X_m|)\) 一致。
于是得到所示恒等式。证毕。
\end{proof}

\begin{theorem}[轨道谱的完全可识别性：由 \texorpdfstring{$(\mathrm{Orb}_m(q))_{q\ge 0}$}{(Orb_m(q))} 递归恢复退化直方图]\label{thm:op-algebra-fold-orbit-spectrum-identifiability-histogram}
对 \(d\ge 1\) 定义退化直方图
\[
N_d(m):=\#\{x\in X_m:\ d_m(x)=d\}.
\]
令
\[
F_m(z):=\sum_{q\ge 0}\mathrm{Orb}_m(q)\frac{z^q}{q!},
\qquad
H_m(y):=F_m(\log(1+y))\in\QQ[[y]].
\]
并定义截断指数多项式
\[
E_d(y):=\sum_{k=0}^{d}\frac{y^k}{k!}.
\]
则有严格因式分解
\[
\boxed{\;
H_m(y)=\prod_{d\ge 1}\bigl(E_d(y)\bigr)^{N_d(m)}\; }.
\]
并且序列 \((N_d(m))_{d\ge 1}\) 由轨道谱 \((\mathrm{Orb}_m(q))_{q\ge 0}\) 唯一确定。
更具体地，存在三角递归：对每个 \(d\ge 1\)，令
\[
R_d(y):=\log H_m(y)-\sum_{j=1}^{d-1}N_j(m)\log E_j(y),
\]
则
\[
\boxed{\;
N_d(m)=-(d+1)!\,[y^{d+1}]\,R_d(y)\; }.
\]
因此 \(N_d(m)\) 仅依赖于有限数据 \(\mathrm{Orb}_m(0),\dots,\mathrm{Orb}_m(d+1)\)。
\end{theorem}
\begin{proof}
由定理 \ref{thm:op-algebra-fold-orbit-mgf-fiber-factorization}，取 \(y=e^{z}-1\) 得
\(
F_m(z)=\prod_{x\in X_m}E_{d_m(x)}(y)
\)；
再令 \(z=\log(1+y)\) 即得
\(
H_m(y)=\prod_{x\in X_m}E_{d_m(x)}(y)=\prod_d E_d(y)^{N_d(m)}
\)。
取对数得到
\(
\log H_m(y)=\sum_d N_d(m)\log E_d(y)
\)。
\par
注意到
\(
E_d(y)=e^{y}-\sum_{k\ge d+1}y^k/k!
=e^{y}-y^{d+1}/(d+1)!+O(y^{d+2})
\)，
从而
\(
\log E_d(y)=y-\frac{y^{d+1}}{(d+1)!}+O(y^{d+2})
\)。
因此对 \(j>d\)，\(\log E_j(y)=y+O(y^{d+2})\)，其在 \(y^{d+1}\) 阶无贡献。
而
\(
R_d(y)=\sum_{j\ge d}N_j(m)\log E_j(y)
\)
的 \(y^{d+1}\) 系数仅来自 \(N_d(m)\log E_d(y)\)，并等于 \(-N_d(m)/(d+1)!\)。
于是 \(N_d(m)=-(d+1)![y^{d+1}]R_d(y)\)。
最后，\([y^{d+1}]\log H_m(y)\) 仅依赖于 \(H_m\) 的前 \(d+1\) 阶系数，而 \(H_m(y)=F_m(\log(1+y))\) 的这些系数仅依赖于 \(F_m\) 的前 \(d+1\) 阶系数，即 \(\mathrm{Orb}_m(0),\dots,\mathrm{Orb}_m(d+1)\)。证毕。
\end{proof}
\leanverified{paper\_op\_algebra\_fold\_orbit\_spectrum\_identifiability\_histogram}

\begin{corollary}[轨道谱诱导的完备离散证书：规范群体积与碰撞矩可重构]\label{cor:op-algebra-fold-orbit-spectrum-reconstructs-gauge-volume-moments}
在定理 \ref{thm:op-algebra-fold-orbit-spectrum-identifiability-histogram} 的记号下，轨道谱 \((\mathrm{Orb}_m(q))_{q\ge 0}\) 唯一确定：
\begin{enumerate}
  \item \textbf{规范群体积：}
  \[
  |G_m|=\prod_{x\in X_m}d_m(x)!\;=\;\prod_{d\ge 1}(d!)^{N_d(m)}.
  \]
  \item \textbf{全阶碰撞矩（功率和）：}对任意整数 \(r\ge 1\)，
  \[
  S_r(m):=\sum_{x\in X_m}d_m(x)^r=\sum_{d\ge 1} d^{\,r}\,N_d(m).
  \]
\end{enumerate}
\end{corollary}
\begin{proof}
由定理 \ref{thm:op-algebra-fold-orbit-spectrum-identifiability-histogram} 可恢复全部 \(N_d(m)\)，代入定义即得。证毕。
\end{proof}
\leanverified{paper\_op\_algebra\_fold\_orbit\_spectrum\_reconstructs\_gauge\_volume\_moments}

\begin{proposition}[轨道数列的 Stieltjes 矩结构：Hankel 半正定与对数凸性]\label{prop:op-algebra-fold-orbit-stieltjes-hankel-logconvex}
序列 \((\mathrm{Orb}_m(q))_{q\ge 0}\) 为 Stieltjes 矩序列：存在取值于 \(\NN\) 的非负随机变量 \(F\) 使得 \(\mathrm{Orb}_m(q)=\mathbb E[F^q]\)。
因此对任意 \(n\ge 0\)，Hankel 矩阵 \(\bigl(\mathrm{Orb}_m(i+j)\bigr)_{0\le i,j\le n}\) 半正定，并且对任意 \(q\ge 1\) 有对数凸不等式
\[
\mathrm{Orb}_m(q)^2\le \mathrm{Orb}_m(q-1)\,\mathrm{Orb}_m(q+1).
\]
\end{proposition}
\begin{proof}
由定理 \ref{thm:op-algebra-fold-orbit-mgf-fiber-factorization} 可取 \(F:=F(\sigma)=|\mathrm{Fix}_{\Omega_m}(\sigma)|\)。
对任意复系数 \((a_i)_{i=0}^n\) 有
\[
\sum_{i,j=0}^{n}a_i\overline{a_j}\,\mathrm{Orb}_m(i+j)
=\mathbb E\!\left[\left|\sum_{i=0}^{n}a_i F^i\right|^2\right]\ge 0,
\]
从而 Hankel 半正定。
对数凸性由 Cauchy--Schwarz
\(
\mathbb E[F^q]^2\le \mathbb E[F^{q-1}]\,\mathbb E[F^{q+1}]
\)
得到。证毕。
\end{proof}
\leanverified{paper\_op\_algebra\_fold\_orbit\_stieltjes\_hankel\_logconvex}

\begin{conjecture}[退化谱的渐近近正则性：最小退化与平均退化同阶]\label{conj:op-algebra-fold-orbit-asymptotic-near-regularity}
在本文的 Fibonacci 稳定型折叠族（包含 bin-fold）中，记平均退化
\(
\bar d_m:=2^m/|X_m|
\)
与最小退化 \(d_{\min}(m)\)。
猜想存在常数 \(C>0\) 使得
\[
\frac{d_{\min}(m)}{\bar d_m}\longrightarrow 1,
\qquad
\max_{x\in X_m}\bigl|d_m(x)-\bar d_m\bigr|\le C,
\qquad(m\to\infty).
\]
在该假设下，对任意固定 \(q\) 充分大 \(m\) 时，低阶轨道谱严格落入定理 \ref{thm:op-algebra-fold-orbit-touchard-loworder} 的 Touchard 通用区间；
并且退化直方图的“可见区”必须随 \(m\) 推进到 \(q>d_{\min}(m)\) 之后才被触发（命题 \ref{prop:op-algebra-fold-orbit-first-deviation-minsector}）。
\end{conjecture}
